% proofs/03_caught.tex
% Who gets caught: signing the composition ratio of the window margin.
% Fragment, to be \input by appendix.tex.
%
% appendix.tex must supply: amsmath, amssymb, amsthm (the proof environment), enumitem,
%   \newtheorem{corollary}{Corollary}, and the macros \Eop, \Prb, \ind.
% Cross-references to the partition result, the clock equivalence and the clock
%   theorem are written in words here.
% Notation notes: (a) B below is the set of newly caught histories. It
%   is unrelated to the stake symbols B_j(s,t) and B^F. (b) The caught share is
%   written \varphi, not \beta, because \beta is the signal weight
%   sigma_v^2/(sigma_v^2 + sigma_eps^2) in the model section. Rename either if the
%   collision reads badly.

\subsection{Who gets caught}
\label{sec:caught}

The clock theorem makes window attenuation equivalent to $W_T C_T \le 1$ and signs only the weight
leg. This section signs the composition leg. It writes the longer clock's pooled sensitivity as a
mass-weighted average of what the shorter clock leaves in the pool and what it takes out, and reads
both $C_T \le 1$ and $W_T C_T \le 1$ off that average.

\paragraph{Setup.}
Fix the threshold margin $\tau$ and the plan and cutoff policies, and compare two window margins
$T' < T$ at a common $\kappa$. Write
\begin{equation}
\label{eq:caught-sets}
A \;=\; \mathcal{C}_P(\tau,T),
\qquad
B \;=\; \mathcal{C}_F(\tau,T')\setminus\mathcal{C}_F(\tau,T),
\end{equation}
so $A$ is the pooled cell under the longer clock and $B$ collects the histories the shorter clock
newly flags. Write
\begin{equation}
\label{eq:caught-phi}
\varphi \;=\; \frac{\Prb(B)}{\Prb(A)}
\end{equation}
for the share of the longer clock's pooled cell that the shorter clock catches.

Each rule pins its own premium kernel. The control-node information set carries the coordinate
``the filing has landed by date $d$'', so the pooled history at window margin $T$ and the pooled
history at window margin $T'$ are different observables, and the engagement posterior
$\pi = \Prb(a = 1 \mid \cdot)$ and the price are taken with respect to whichever of the two the
rule supplies. Write $h^{(T)}$ and $h^{(T')}$ for the kernel $h = \pi p$ evaluated at the
control-node information set of the rule $(\tau,T)$ and of the rule $(\tau,T')$. The two pooled
averages of the clock theorem are then
\begin{equation}
\label{eq:caught-MP}
M_P(\tau,T) \;=\; \Delta_m\,\Eop\bigl[h^{(T)} \bigm| A\bigr],
\qquad
M_P(\tau,T') \;=\; \Delta_m\,\Eop\bigl[h^{(T')} \bigm| A\setminus B\bigr],
\end{equation}
the second reading of the pooled cell at the shorter clock being part~(i) below. Set
\begin{equation}
\label{eq:caught-sens}
s_A \;=\; \partial_\kappa\Eop\bigl[h^{(T)} \bigm| A\bigr],
\qquad
s_{A\setminus B} \;=\; \partial_\kappa\Eop\bigl[h^{(T')} \bigm| A\setminus B\bigr],
\end{equation}
so that, with $\Delta_m > 0$ finite,
\begin{equation}
\label{eq:caught-SP}
\mathcal{S}_P(\tau,T) \;=\; \Delta_m\,|s_A|,
\qquad
\mathcal{S}_P(\tau,T') \;=\; \Delta_m\,|s_{A\setminus B}|,
\qquad
C_T \;=\; \frac{|s_{A\setminus B}|}{|s_A|} ,
\end{equation}
by the definitions of the two pooled sensitivities and no more.

Two further quantities separate what the catch removes from how the surviving pool is re-read:
\begin{equation}
\label{eq:caught-raw}
\tilde s_B \;=\; \partial_\kappa\Eop\bigl[h^{(T)} \bigm| B\bigr],
\qquad
\delta \;=\; \partial_\kappa\Eop\bigl[h^{(T')} - h^{(T)} \bigm| A\setminus B\bigr] .
\end{equation}
Here $\tilde s_B$ is the \emph{caught-only leg}, the noise sensitivity of the newly caught
histories priced as the longer clock priced them, and $\delta$ is the \emph{re-pricing term}: the pooled posterior conditions
on the pool the rule leaves behind, so the histories that stay pooled are read against a smaller
pool once the clock shortens, and $\delta$ is the rate at which that re-reading moves their noise
sensitivity. Define the \emph{net cut leg} by
\begin{equation}
\label{eq:caught-sB}
s_B \;:=\; \frac{1}{\Prb(B)}\,
\partial_\kappa\Bigl(
\Eop\bigl[h^{(T)}\ind_A\bigr] - \Eop\bigl[h^{(T')}\ind_{A\setminus B}\bigr]
\Bigr) .
\end{equation}
The pooled cell contributes to the expected engagement premium an amount
\[
\begin{aligned}
\Delta_m\,\Eop\bigl[h^{(T)}\ind_A\bigr]
  \;&=\; \bigl(1-\Omega(\tau,T)\bigr)\,M_P(\tau,T)
  &&\text{at the longer clock,} \\[3pt]
\Delta_m\,\Eop\bigl[h^{(T')}\ind_{A\setminus B}\bigr]
  \;&=\; \bigl(1-\Omega(\tau,T')\bigr)\,M_P(\tau,T')
  &&\text{at the shorter one.}
\end{aligned}
\]
So the net cut leg $s_B$ is the derivative of $\Lambda_T - \Lambda_{T'}$ per unit of pooled mass removed: it is
the noise sensitivity of the premium mass the shorter clock removes from the pool, and it
carries the survivors' re-pricing. It is not the caught-only leg $\tilde s_B$. Under the extra clause of (C-2) it
splits into the two parts just named,
\begin{equation}
\label{eq:caught-split}
s_B \;=\; \tilde s_B \;-\; \frac{1-\varphi}{\varphi}\,\delta ,
\end{equation}
which part~(iii) proves.

The hypotheses are:
\begin{enumerate}[label=(C-\arabic*),leftmargin=3.2em,itemsep=3pt]
\item \emph{Fixed policies.} The plan menu, the execution policies and the cutoff vector are held
  fixed in $\kappa$ and held fixed across the two window margins, the threshold margin $\tau$ is
  common, the noise intensity $\kappa$ is common, and the timing carries no feedback from order
  flow into the blockholder's executed order path. Liquidity enters the primitives in one place
  only, the law of the ternary noise mark, which is standing condition (S10) of the partition
  subsection; fixed policies are standing condition (S11), cited separately here. Part~(ii)
  consumes both conditions.
\item \emph{Differentiability.} The two conditional expectations in \eqref{eq:caught-sens} are
  differentiable in $\kappa$ at the point compared. For the split \eqref{eq:caught-split} into
  $\tilde s_B$ and $\delta$, $\Eop[h^{(T)}\mid B]$ is also differentiable. The kernel is the
  pinned version of
  standing condition (S8), so $h^{(T)}$ and $h^{(T')}$ are well-defined random variables rather
  than equivalence classes; part~(iii) consumes that.
\item \emph{Non-degeneracy.} $\Prb(B) > 0$, $0 < \Omega(\tau,T)$, $\Omega(\tau,T') < 1$, and
  $s_A \ne 0$, equivalently $\mathcal{S}_P(\tau,T) > 0$. The clause $0 < \Omega(\tau,T)$ is there
  so that the clock theorem applies. Wherever this corollary is used to read the criterion
  $W_T C_T \le 1$ in Theorem~\ref{thm:clock}, that theorem's hypotheses also apply. They include
  the hypotheses of Proposition~\ref{prop:factorisation} at both clocks, including the invariance
  of the flagged endpoint in $\kappa$ at fixed policies.
\end{enumerate}

\begin{corollary}[Who gets caught]
\label{cor:caught}
Under (C-1) to (C-3):
\begin{enumerate}[label=(\roman*),leftmargin=2.4em,itemsep=3pt]
\item \emph{The cut is nested.} $\mathcal{C}_F(\tau,T)\subseteq\mathcal{C}_F(\tau,T')$, hence
  $B \subseteq A$ and $A\setminus B = \mathcal{C}_P(\tau,T')$, and the weight leg is
  \begin{equation}
  \label{eq:caught-WT}
  W_T \;=\; 1-\varphi \;\le\; 1 ,
  \qquad \varphi\in(0,1).
  \end{equation}
\item \emph{The cell masses are $\kappa$-free.}
  $\partial_\kappa\Prb(A) = \partial_\kappa\Prb(B) = 0$, so $\varphi$ does not depend on $\kappa$.
\item \emph{The cut identity.}
  \begin{equation}
  \label{eq:caught-identity}
  s_{A\setminus B}
  \;=\;
  \frac{\Prb(A)\,s_A - \Prb(B)\,s_B}{\Prb(A)-\Prb(B)}
  \;=\;
  \frac{s_A - \varphi\,s_B}{1-\varphi}
  \;=\;
  \frac{s_A - \varphi\,\tilde s_B}{1-\varphi} \;+\; \delta ,
  \end{equation}
  the last expression under the extra differentiability of (C-2). Equivalently
  $s_A = (1-\varphi)\,s_{A\setminus B} + \varphi\,s_B$: the longer clock's pooled sensitivity is
  the mass-weighted average of the sensitivity of what the shorter clock leaves in the pool and
  the net cut leg $s_B$, the derivative of $\Lambda_T-\Lambda_{T'}$ per unit of pooled mass removed, which
  carries the survivors' re-pricing.
\item \emph{The composition ratio.}
  \begin{equation}
  \label{eq:caught-CT}
  \begin{aligned}
  C_T \;&=\; \frac{|s_A-\varphi s_B|}{(1-\varphi)\,|s_A|}, \\[3pt]
  C_T \le 1 \quad&\Longleftrightarrow\quad
  (s_B-s_A)\bigl(\varphi s_B-(2-\varphi)s_A\bigr) \le 0 ,
  \end{aligned}
  \end{equation}
  that is, if and only if $s_B$ lies weakly between $s_A$ and
  $\bigl((2-\varphi)/\varphi\bigr)s_A$.
\item \emph{The attenuation criterion.}
  \begin{equation}
  \label{eq:caught-WTCT}
  \begin{aligned}
  W_T C_T \;&=\; \frac{|s_A-\varphi s_B|}{|s_A|}, \\[3pt]
  W_T C_T \le 1 \quad&\Longleftrightarrow\quad
  s_B\bigl(2s_A-\varphi s_B\bigr) \ge 0 ,
  \end{aligned}
  \end{equation}
  that is, if and only if $s_B$ lies weakly between $0$ and $(2/\varphi)s_A$.
\item \emph{Readings.} Write $\rho = s_B/s_A$, which is well defined by (C-3).
  \begin{enumerate}[label=(\alph*),leftmargin=2.2em,itemsep=2pt]
  \item If $s_A s_B \le 0$, including $s_B = 0$, then $C_T > 1$ strictly.
  \item If $s_A s_B > 0$, then
    $C_T \le 1$ if and only if
    $|s_A| \le |s_B| \le \bigl((2-\varphi)/\varphi\bigr)|s_A|$.
  \item If the cut does not reverse the pooled sensitivity, $s_{A\setminus B}\,s_A \ge 0$, then
    $C_T \le 1$ if and only if $\rho \ge 1$. The inequality $\rho \ge 1 > 0$ already forces a
    common sign, so that reading is $|s_B| \ge |s_A|$. Under a common sign the proviso
    $\sigma \ge 0$ is $|s_B| \le |s_A|/\varphi$.
  \item $W_T C_T \le 1$ if and only if $s_B$ shares the sign of $s_A$ or is zero, and
    $|s_B| \le (2/\varphi)|s_A|$.
  \item The two upper limits $(2-\varphi)/\varphi$ and $2/\varphi$ are strictly decreasing in
    $\varphi$ and diverge as $\varphi\downarrow0$. For every $M \ge 1$ and every cut with
    $\varphi \le 2/(1+M)$ and $|s_B| \le M|s_A|$, the upper limits in (b) and (d) are slack, so
    there $C_T \le 1$ reads: a common sign together with $|s_B| \ge |s_A|$; and $W_T C_T \le 1$
    reads: a common sign or $s_B = 0$.
  \end{enumerate}
\end{enumerate}
\end{corollary}

\begin{proof}
\emph{Part (i).} The disclosure indicator at window margin $T$ is
$D(\tau,T) = \ind\{a=1,\ c(\tau)+T\le H\}$, where $c(\tau)$ is the round in which the executed
stake path crosses $\tau$. If $D(\tau,T)=1$ then $c(\tau)+T' \le c(\tau)+T \le H$ because
$T'<T$, so $D(\tau,T')=1$. Hence $\mathcal{C}_F(\tau,T)\subseteq\mathcal{C}_F(\tau,T')$ and, taking
complements, $\mathcal{C}_P(\tau,T')\subseteq\mathcal{C}_P(\tau,T)=A$. Therefore
$B = \mathcal{C}_F(\tau,T')\cap A \subseteq A$ and
$A\setminus B = A\cap\mathcal{C}_P(\tau,T') = \mathcal{C}_P(\tau,T')$. Taking probabilities,
$\Omega(\tau,T) \le \Omega(\tau,T') < 1$, so both pooled cells carry positive mass; since
$B\subseteq A$, $\Prb(A\setminus B) = \Prb(A)-\Prb(B)$, and
\[
W_T \;=\; \frac{1-\Omega(\tau,T')}{1-\Omega(\tau,T)}
\;=\; \frac{\Prb(A\setminus B)}{\Prb(A)}
\;=\; 1-\varphi .
\]
By (C-3), $\Prb(B)>0$ gives $\varphi>0$ and $\Omega(\tau,T')<1$ gives $\Prb(A\setminus B)>0$, hence
$\varphi<1$.

\emph{Part (ii).} Under fixed policies the selected plan and the executed order path are functions
of the plan index and the signal alone, and by the no-feedback clause of (C-1) the order placed in
a round does not read realised order flow. So the crossing round $c(\tau)$, and with it
$D(\tau,T)$ and $D(\tau,T')$, are measurable with respect to the plan and signal. The intensity
$\kappa$ parametrises the noise law only; the laws of the value, the signal noise and the bidder
draw carry no $\kappa$, and (C-1) holds the cutoff vector fixed in $\kappa$, so the joint law of
the plan and the signal is one and the same at every $\kappa$. Hence $\Prb(A)$ and $\Prb(B)$ do not
vary with $\kappa$, and neither does their ratio $\varphi$. This is the same constancy the clock
theorem uses for the flagged weight, applied at both window margins.

\emph{Part (iii).} By part (i), $B$ and $A\setminus B$ partition $A$ and both have positive mass.
Write the pooled contributions to the expected premium at the two clocks, in kernel units, as
\[
\begin{aligned}
\Lambda_T \;&=\; \Eop\bigl[h^{(T)}\ind_A\bigr]
  \;=\; \Prb(A)\,\Eop\bigl[h^{(T)}\mid A\bigr], \\[3pt]
\Lambda_{T'} \;&=\; \Eop\bigl[h^{(T')}\ind_{A\setminus B}\bigr]
  \;=\; \Prb(A\setminus B)\,\Eop\bigl[h^{(T')}\mid A\setminus B\bigr] .
\end{aligned}
\]
By part (ii) the two masses are constants in $\kappa$, so the derivative that (C-2) supplies passes
through them and leaves them alone:
\[
\partial_\kappa \Lambda_T \;=\; \Prb(A)\,s_A ,
\qquad
\partial_\kappa \Lambda_{T'} \;=\; \bigl(\Prb(A)-\Prb(B)\bigr)\,s_{A\setminus B} .
\]
Definition \eqref{eq:caught-sB} reads
$s_B = \partial_\kappa(\Lambda_T - \Lambda_{T'})/\Prb(B)$, a division by a positive number, so
\[
\Prb(B)\,s_B \;=\; \Prb(A)\,s_A - \bigl(\Prb(A)-\Prb(B)\bigr)\,s_{A\setminus B} ,
\]
and solving for $s_{A\setminus B}$, which the positive mass $\Prb(A)-\Prb(B)$ permits, gives the
first equality of \eqref{eq:caught-identity}. Dividing numerator and denominator by $\Prb(A)$ gives
the second. Multiplying through by $1-\varphi$ and collecting terms gives
$s_A = (1-\varphi)s_{A\setminus B}+\varphi s_B$.

For the third expression, add and subtract $\Eop[h^{(T)}\ind_{A\setminus B}]$ inside
$\Lambda_T-\Lambda_{T'}$. Since $\ind_A = \ind_{A\setminus B}+\ind_B$ by part (i),
\[
\begin{aligned}
\Lambda_T - \Lambda_{T'}
\;&=\; \Eop\bigl[h^{(T)}\ind_B\bigr]
  \;-\; \Eop\bigl[\bigl(h^{(T')}-h^{(T)}\bigr)\ind_{A\setminus B}\bigr] \\[3pt]
\;&=\; \Prb(B)\,\Eop\bigl[h^{(T)}\mid B\bigr]
  \;-\; \Prb(A\setminus B)\,\Eop\bigl[h^{(T')}-h^{(T)}\mid A\setminus B\bigr] .
\end{aligned}
\]
Both masses are $\kappa$-free by part (ii), and the extra clause of (C-2) makes the first term
differentiable, hence so is the second. Differentiating and dividing by $\Prb(B)$ gives
$s_B = \tilde s_B - \bigl((1-\varphi)/\varphi\bigr)\delta$, which is \eqref{eq:caught-split}.
Substituting it into the second equality of \eqref{eq:caught-identity} gives
the third.

\emph{Part (iv).} By \eqref{eq:caught-SP} and part (i), $C_T = |s_{A\setminus B}|/|s_A|$, and by
\eqref{eq:caught-identity} with $1-\varphi>0$,
\[
C_T \;=\; \frac{|s_A-\varphi s_B|}{(1-\varphi)|s_A|} .
\]
The denominator is non-zero because $s_A \ne 0$ in (C-3). Both sides of $C_T\le 1$ are
non-negative, so $C_T\le 1$ if and only if $s_{A\setminus B}^2\le s_A^2$, and multiplying by
$(1-\varphi)^2>0$, if and only if $u^2-w^2\le 0$, where $u=s_A-\varphi s_B$ and
$w=(1-\varphi)s_A$. Since $u-w=\varphi(s_A-s_B)$ and $u+w=(2-\varphi)s_A-\varphi s_B$,
\[
u^2-w^2 \;=\; (u-w)(u+w) \;=\; \varphi\,(s_A-s_B)\bigl((2-\varphi)s_A-\varphi s_B\bigr).
\]
Since $\varphi>0$, the sign of the whole is the sign of
$(s_A-s_B)\bigl((2-\varphi)s_A-\varphi s_B\bigr)$, and negating both factors leaves the product
unchanged, which gives the stated form
$(s_B-s_A)\bigl(\varphi s_B-(2-\varphi)s_A\bigr)\le 0$. Writing that product as
$\varphi(s_B-s_A)\bigl(s_B-\bigl((2-\varphi)/\varphi\bigr)s_A\bigr)$ shows it is non-positive
exactly when $s_B$ lies weakly between the two roots $s_A$ and
$\bigl((2-\varphi)/\varphi\bigr)s_A$.

\emph{Part (v).} By part (i) and part (iv),
\[
W_T C_T \;=\; (1-\varphi)\,\frac{|s_{A\setminus B}|}{|s_A|}
\;=\; \frac{\bigl|(1-\varphi)s_{A\setminus B}\bigr|}{|s_A|}
\;=\; \frac{|s_A-\varphi s_B|}{|s_A|} .
\]
As in part (iv), $W_T C_T\le 1$ if and only if $u^2-s_A^2\le 0$ for $u=s_A-\varphi s_B$, and since
$u-s_A=-\varphi s_B$ and $u+s_A=2s_A-\varphi s_B$,
\[
u^2-s_A^2 \;=\; (u-s_A)(u+s_A) \;=\; -\varphi\,s_B\bigl(2s_A-\varphi s_B\bigr),
\]
so the criterion is $s_B(2s_A-\varphi s_B)\ge 0$. Writing it as
$\varphi\,s_B\bigl((2/\varphi)s_A-s_B\bigr)\ge 0$ shows it holds exactly when $s_B$ lies weakly
between $0$ and $(2/\varphi)s_A$.

\emph{Part (vi).} Since $\varphi\in(0,1)$, the factor $(2-\varphi)/\varphi$ is strictly positive,
so the two roots in part (iv) are non-zero and share the sign of $s_A$, and the closed interval
they bound does not contain zero. If $s_As_B\le 0$ then $s_B$ is zero or lies on the far side of
zero from both roots, so it is strictly outside that interval, the product in
\eqref{eq:caught-CT} is strictly positive, and $C_T>1$. That is (a). For (b), divide the interval
condition by $s_A$: with a common sign, $\rho>0$ and $s_B$ between $s_A$ and
$\bigl((2-\varphi)/\varphi\bigr)s_A$ reads $1\le\rho\le(2-\varphi)/\varphi$, and multiplying by
$|s_A|$ gives the stated magnitudes. For (c), set
$\sigma = s_{A\setminus B}/s_A = (1-\varphi\rho)/(1-\varphi)$ by \eqref{eq:caught-identity}. The
proviso $s_{A\setminus B}s_A\ge0$ is $\sigma\ge0$, so $C_T = |\sigma| = \sigma$, and $C_T\le1$ if
and only if $1-\varphi\rho\le1-\varphi$, that is $\rho\ge1$. The inequality $\rho\ge1>0$ already
forces a common sign, under which $\rho\ge1$ is $|s_B|\ge|s_A|$. Under a common sign the proviso
$\sigma\ge0$ is $\varphi\rho\le1$, that is $|s_B|\le|s_A|/\varphi$. For (d), divide the interval
condition of part (v) by $s_A$: it reads $0\le\rho\le2/\varphi$, which is a common sign or
$s_B=0$, together with $|s_B|\le(2/\varphi)|s_A|$. For (e), both $(2-\varphi)/\varphi$ and
$2/\varphi$ are strictly decreasing on $(0,1)$ and tend to $+\infty$ as $\varphi\downarrow0$. Fix
$M\ge1$ and $\varphi\le2/(1+M)$. Then $(2-\varphi)/\varphi\ge M$, because that inequality is
$2-\varphi\ge M\varphi$, and $2/\varphi\ge(2-\varphi)/\varphi\ge M$. So $|s_B|\le M|s_A|$ makes
the upper limits in (b) and (d) slack, leaving in (b) the common sign with $|s_B|\ge|s_A|$ and in
(d) the common sign or $s_B=0$.
\end{proof}

\paragraph{Reading.}
Part~(iii) is the whole content in one line: the longer clock's pooled sensitivity is an average
over what the shorter clock leaves in the pool and what it takes out of the pool, with weights the
two masses, and those weights do not move with liquidity. So shortening the clock pulls the pooled
sensitivity down only if what it takes out sat above the average. Parts~(iv) and~(vi) put the
bound on how far above. What the clock removes must move with the pool in $\kappa$ and be at least
as noise-sensitive as the pool, and it must not be so much more sensitive that the remaining pool
is driven through zero to a larger magnitude than the pool started with. A cut that takes a small
share of the pool faces only the first two requirements when, as part~(vi)(e) states, there is
some $M\ge 1$ with $\varphi \le 2/(1+M)$ and $|s_B|\le M|s_A|$. Part~(v) does the same for the criterion
the clock theorem states, $W_T C_T\le1$, where the weight leg is already working in the same
direction: there what the clock removes need only move with the pool and stay inside $2/\varphi$
times its sensitivity.

What the clock removes has two parts, and \eqref{eq:caught-split} separates them. The first is the
caught-only leg, the noise sensitivity of the newly caught histories read at the prices the longer
clock gave them. The second is the re-pricing term: the surviving pool is a smaller pool once the
clock shortens, and the market reads it against that smaller pool, so its own sensitivity moves
too. The criterion runs on the sum, which is the sensitivity of the premium mass that leaves the
pool. When the re-pricing term is zero the sum is the first part alone, and the criterion is
a statement about the caught histories and nothing else.

The corollary turns the ratio $C_T$ into a question with an answer at every calibration node, and
that question is who gets caught. Under a common sign, $s_A s_B > 0$, the composition ratio
satisfies $C_T \le 1$ if and only if the net cut leg $s_B$ lies weakly between $s_A$ and
$\bigl((2-\varphi)/\varphi\bigr)\,s_A$, which is part~(iv); and the shorter clock lowers noise
sensitivity overall, $W_T C_T \le 1$, if and only if $s_B$ lies weakly between $0$ and
$(2/\varphi)\,s_A$, which is part~(v). The two bands differ: the weight leg lets the overall
criterion accept a net cut leg less sensitive than the pool, which the composition criterion alone
does not.

% =============================================================================
% The same accounting law at the threshold margin.
% =============================================================================

\csname unlabelledtrue\endcsname
\begin{corollary}[Who gets caught at the threshold margin]
\label{cor:caught-threshold}
Fix a common window margin $T$ and liquidity $\kappa$, and compare $b_0<\tau'<\tau$, with
$(\tau',T)$ the tighter rule. Write
\begin{equation}
\label{eq:caught-tau-sets}
A_\tau=\mathcal C_P(\tau,T),
\qquad
B_\tau=\mathcal C_F(\tau',T)\setminus\mathcal C_F(\tau,T),
\qquad
R_\tau=A_\tau\setminus B_\tau,
\end{equation}
and let $h^{(\tau)}$ and $h^{(\tau')}$ be the kernels supplied by the two rules. Set
\begin{align}
\label{eq:caught-tau-definitions}
\varphi_\tau&=\frac{\Prb(B_\tau)}{\Prb(A_\tau)},
&s_{A,\tau}&=\partial_\kappa\Eop[h^{(\tau)}\mid A_\tau],
&s_{R,\tau}&=\partial_\kappa\Eop[h^{(\tau')}\mid R_\tau],
\nonumber\\
\widetilde s_{B,\tau}&=\partial_\kappa\Eop[h^{(\tau)}\mid B_\tau],
&\delta_\tau&=\partial_\kappa\Eop[h^{(\tau')}-h^{(\tau)}\mid R_\tau],
\end{align}
and define the net cut leg by
\begin{equation}
\label{eq:caught-tau-sB}
s_{B,\tau}:=\frac{1}{\Prb(B_\tau)}\partial_\kappa\left(
\Eop[h^{(\tau)}\ind_{A_\tau}]
-\Eop[h^{(\tau')}\ind_{R_\tau}]
\right).
\end{equation}
Assume standing conditions (S1) to (S11) and the following conditions.
\begin{enumerate}[label=(C$\tau$-\arabic*),leftmargin=3.8em,itemsep=3pt]
\item \emph{A fixed threshold cut.} The window margin and liquidity are common, and
  $b_0<\tau'<\tau$. Conditions (S4), (S5), and (S7) apply at both thresholds, so both disclosure
  indicators are functions of the same plan and signal. Liquidity enters only through the noise law
  by (S10), and the policies are common by (S11).
\item \emph{Differentiability and pinned kernels.} The maps
  $\kappa\mapsto\Eop[h^{(\tau)}\mid A_\tau]$ and
  $\kappa\mapsto\Eop[h^{(\tau')}\mid R_\tau]$ are differentiable at the compared liquidity.
  Condition (S8) pins both kernels. For the split into the caught-only leg and the re-pricing term,
  $\Eop[h^{(\tau)}\mid B_\tau]$ is also differentiable.
\item \emph{Non-degeneracy and factorisation.}
  $0<\Prb(B_\tau)<\Prb(A_\tau)$, $0<\Omega(\tau,T)$,
  $\Omega(\tau',T)<1$, and $s_{A,\tau}\ne0$. Whenever the aggregate criterion is invoked, the
  hypotheses of Proposition~\ref{prop:factorisation}, including flagged-endpoint invariance, hold
  at both thresholds.
\item \emph{Erasure representation.} Only the identification with
  Condition~\ref{cond:D} requires standing conditions (S12) and (S13), order size two, and
  $\kappa\in(0,1)$, so that Lemmas~\ref{lem:g1} to~\ref{lem:g3} apply. The cut algebra below does
  not use this clause.
\end{enumerate}
Under (C$\tau$-1) to (C$\tau$-3):
\begin{enumerate}[label=(\roman*),leftmargin=2.4em,itemsep=3pt]
\item \emph{The cut is nested.}
  $\mathcal C_F(\tau,T)\subseteq\mathcal C_F(\tau',T)$, hence
  $B_\tau\subseteq A_\tau$ and $R_\tau=\mathcal C_P(\tau',T)$, and
  \begin{equation}
  \label{eq:caught-tau-W}
  W_\tau=\frac{1-\Omega(\tau',T)}{1-\Omega(\tau,T)}
  =1-\varphi_\tau\le1,
  \qquad \varphi_\tau\in(0,1).
  \end{equation}
\item \emph{The cell masses are $\kappa$-free.}
  $\partial_\kappa\Prb(A_\tau)=\partial_\kappa\Prb(B_\tau)=0$, so
  $\varphi_\tau$ does not depend on $\kappa$.
\item \emph{The cut identity.}
  \begin{equation}
  \label{eq:caught-tau-identity}
  s_{R,\tau}
  =\frac{\Prb(A_\tau)s_{A,\tau}-\Prb(B_\tau)s_{B,\tau}}
  {\Prb(A_\tau)-\Prb(B_\tau)}
  =\frac{s_{A,\tau}-\varphi_\tau s_{B,\tau}}{1-\varphi_\tau}
  =\frac{s_{A,\tau}-\varphi_\tau\widetilde s_{B,\tau}}{1-\varphi_\tau}
  +\delta_\tau,
  \end{equation}
  where the last expression uses the extra differentiability in (C$\tau$-2). Equivalently,
  \begin{equation}
  \label{eq:caught-tau-average}
  s_{A,\tau}=(1-\varphi_\tau)s_{R,\tau}+\varphi_\tau s_{B,\tau},
  \qquad
  s_{B,\tau}=\widetilde s_{B,\tau}
  -\frac{1-\varphi_\tau}{\varphi_\tau}\delta_\tau.
  \end{equation}
\item \emph{The composition ratio.}
  \begin{equation}
  \label{eq:caught-tau-C}
  \begin{aligned}
  C_\tau&=\frac{|s_{A,\tau}-\varphi_\tau s_{B,\tau}|}
  {(1-\varphi_\tau)|s_{A,\tau}|},\\
  C_\tau\le1
  &\quad\Longleftrightarrow\quad
  (s_{B,\tau}-s_{A,\tau})
  \bigl(\varphi_\tau s_{B,\tau}-(2-\varphi_\tau)s_{A,\tau}\bigr)\le0.
  \end{aligned}
  \end{equation}
  Thus $C_\tau\le1$ exactly when the net cut leg lies weakly between $s_{A,\tau}$ and
  $((2-\varphi_\tau)/\varphi_\tau)s_{A,\tau}$. Under (C$\tau$-4), this is equivalent to
  Condition~\ref{cond:D}.
\item \emph{The attenuation criterion.}
  \begin{equation}
  \label{eq:caught-tau-WC}
  \begin{aligned}
  W_\tau C_\tau
  &=\frac{|s_{A,\tau}-\varphi_\tau s_{B,\tau}|}{|s_{A,\tau}|},\\
  W_\tau C_\tau\le1
  &\quad\Longleftrightarrow\quad
  s_{B,\tau}(2s_{A,\tau}-\varphi_\tau s_{B,\tau})\ge0.
  \end{aligned}
  \end{equation}
  Thus attenuation holds exactly when the net cut leg lies weakly between $0$ and
  $(2/\varphi_\tau)s_{A,\tau}$. Under (C$\tau$-3), this is equivalent to the tighter threshold
  weakly lowering aggregate liquidity sensitivity.
\item \emph{Readings.} Write $\rho_\tau=s_{B,\tau}/s_{A,\tau}$, which is well defined by
  (C$\tau$-3).
  \begin{enumerate}[label=(\alph*),leftmargin=2.2em,itemsep=2pt]
  \item If $s_{A,\tau}s_{B,\tau}\le0$, including $s_{B,\tau}=0$, then $C_\tau>1$.
  \item If $s_{A,\tau}s_{B,\tau}>0$, then $C_\tau\le1$ if and only if
    $|s_{A,\tau}|\le|s_{B,\tau}|\le
    ((2-\varphi_\tau)/\varphi_\tau)|s_{A,\tau}|$.
  \item If $s_{R,\tau}s_{A,\tau}\ge0$, then $C_\tau\le1$ if and only if
    $\rho_\tau\ge1$. That inequality forces $s_{B,\tau}$ and $s_{A,\tau}$ to have a common sign and
    says $|s_{B,\tau}|\ge|s_{A,\tau}|$. Under that common sign, the proviso
    $s_{R,\tau}s_{A,\tau}\ge0$ is
    $|s_{B,\tau}|\le|s_{A,\tau}|/\varphi_\tau$.
  \item $W_\tau C_\tau\le1$ if and only if $s_{B,\tau}$ has the sign of
    $s_{A,\tau}$ or is zero, and
    $|s_{B,\tau}|\le(2/\varphi_\tau)|s_{A,\tau}|$.
  \item The upper limits $(2-\varphi_\tau)/\varphi_\tau$ and $2/\varphi_\tau$ fall as
    $\varphi_\tau$ rises and diverge as $\varphi_\tau\downarrow0$. If $M\ge1$,
    $\varphi_\tau\le2/(1+M)$, and $|s_{B,\tau}|\le M|s_{A,\tau}|$, the upper limits are slack.
    The composition condition then asks for a common sign and
    $|s_{B,\tau}|\ge|s_{A,\tau}|$. The aggregate condition asks only for a common sign or
    $s_{B,\tau}=0$.
  \end{enumerate}
\end{enumerate}
\end{corollary}
\csname unlabelledfalse\endcsname

\begin{proof}
By (C$\tau$-1) and Lemma~\ref{lem:threshold-weight}, at either threshold the disclosure indicator
has the product form
\[
D(\tau_0,T)=\ind\{a=1\}\ind\{B(s,H-T)\ge\tau_0\},
\qquad \tau_0\in\{\tau',\tau\}.
\]
If a history is flagged at $\tau$, then its common stake path satisfies
$B(s,H-T)\ge\tau>\tau'$, so it is flagged at $\tau'$. This proves the flagged-set inclusion.
Taking complements gives $R_\tau=\mathcal C_P(\tau',T)$. The mass identity then gives
\eqref{eq:caught-tau-W}, and (C$\tau$-3) gives $0<\varphi_\tau<1$.

Conditions (S3), (S7), (S10), and (S11) make $A_\tau$ and $B_\tau$ events determined by the common plan
and signal, whose law does not move with $\kappa$. Their masses and their ratio are therefore
$\kappa$-free. This proves part~(ii).

For part~(iii), put
\[
\Lambda_\tau=\Eop[h^{(\tau)}\ind_{A_\tau}],
\qquad
\Lambda_{\tau'}=\Eop[h^{(\tau')}\ind_{R_\tau}].
\]
Part~(ii) and (C$\tau$-2) give
\[
\partial_\kappa\Lambda_\tau=\Prb(A_\tau)s_{A,\tau},
\qquad
\partial_\kappa\Lambda_{\tau'}=\Prb(R_\tau)s_{R,\tau}.
\]
Definition~\eqref{eq:caught-tau-sB} and
$\Prb(R_\tau)=\Prb(A_\tau)-\Prb(B_\tau)$ now give the first two equalities in
\eqref{eq:caught-tau-identity} and the average identity in
\eqref{eq:caught-tau-average}. For the split, use
$\ind_{A_\tau}=\ind_{B_\tau}+\ind_{R_\tau}$ to write
\[
\Lambda_\tau-\Lambda_{\tau'}
=\Prb(B_\tau)\Eop[h^{(\tau)}\mid B_\tau]
-\Prb(R_\tau)\Eop[h^{(\tau')}-h^{(\tau)}\mid R_\tau].
\]
Differentiate and divide by $\Prb(B_\tau)$. This proves the second identity in
\eqref{eq:caught-tau-average}; substituting it into the average identity proves the last equality in
\eqref{eq:caught-tau-identity}.

By (S9) and Lemma~\ref{lem:two-cell},
\[
\mathcal S_P(\kappa,\tau,T)=\Delta_m|s_{A,\tau}|,
\qquad
\mathcal S_P(\kappa,\tau',T)=\Delta_m|s_{R,\tau}|.
\]
Use \eqref{eq:caught-tau-identity} and $1-\varphi_\tau>0$ to obtain the first line of
\eqref{eq:caught-tau-C}. To test $C_\tau\le1$, set
$u=s_{A,\tau}-\varphi_\tau s_{B,\tau}$ and
$w=(1-\varphi_\tau)s_{A,\tau}$. Then
\[
u^2-w^2
=\varphi_\tau(s_{A,\tau}-s_{B,\tau})
\bigl((2-\varphi_\tau)s_{A,\tau}-\varphi_\tau s_{B,\tau}\bigr).
\]
The denominator in $C_\tau$ is positive. Squaring therefore preserves the comparison, and the
display is nonpositive exactly when $s_{B,\tau}$ lies between the two roots in part~(iv). Under
(C$\tau$-4), Lemma~\ref{lem:g3} identifies the two pooled sensitivities with the absolute revelation
values, so this is Condition~\ref{cond:D}.

Multiplying the first line of \eqref{eq:caught-tau-C} by
$W_\tau=1-\varphi_\tau$ proves the first line of \eqref{eq:caught-tau-WC}. The second follows from
\[
(s_{A,\tau}-\varphi_\tau s_{B,\tau})^2-s_{A,\tau}^2
=-\varphi_\tau s_{B,\tau}(2s_{A,\tau}-\varphi_\tau s_{B,\tau}).
\]
The two roots are $0$ and $(2/\varphi_\tau)s_{A,\tau}$. The factorisation assumed in
(C$\tau$-3) turns $W_\tau C_\tau$ into the ratio of aggregate sensitivities, proving part~(v).

For part~(vi), the two roots in part~(iv) are nonzero and have the sign of $s_{A,\tau}$, so zero
lies strictly outside the closed interval between them. This proves (a). Under a common sign,
division by $s_{A,\tau}$ turns their interval into
$1\le\rho_\tau\le(2-\varphi_\tau)/\varphi_\tau$, proving (b). The average identity gives
\[
\frac{s_{R,\tau}}{s_{A,\tau}}
=\frac{1-\varphi_\tau\rho_\tau}{1-\varphi_\tau}.
\]
Under the proviso in (c), the left side is nonnegative, so its absolute value is at most one exactly
when $\rho_\tau\ge1$; the same display gives the stated upper limit. Dividing the interval in
part~(v) by $s_{A,\tau}$ proves (d). Finally, both upper limits fall in $\varphi_\tau$ and diverge at
zero. The inequality $\varphi_\tau\le2/(1+M)$ implies
$(2-\varphi_\tau)/\varphi_\tau\ge M$ and $2/\varphi_\tau\ge M$, which proves (e).
\end{proof}

% =============================================================================
% The split and the polynomial implications for either dial.
% =============================================================================

\csname unlabelledtrue\endcsname
\begin{lemma}[The cut split at both margins]
\label{lem:cut-split}
Assume standing conditions (S1) to (S11). Compare either a shorter clock $T'<T$ at common
$(\tau,\kappa)$ or a tighter threshold $b_0<\tau'<\tau$ at common $(T,\kappa)$. Call the looser rule
$r_-$ and the tighter rule $r_+$, and write
\[
A=\mathcal C_{P,r_-},
\qquad
B=\mathcal C_{F,r_+}\setminus\mathcal C_{F,r_-},
\qquad
R=A\setminus B,
\qquad
\varphi=\frac{\Prb(B)}{\Prb(A)}.
\]
Here $R$ is the surviving tighter pool; Lemma~\ref{lem:silence} uses $A$ for that role.
Let $h^-$ and $h^+$ be the kernels supplied by $r_-$ and $r_+$. The following conditions are in
force.
\begin{enumerate}[label=(K-\arabic*),leftmargin=3.2em,itemsep=3pt]
\item \emph{A non-null nested cut.} Conditions (S4), (S5), (S7), and (S11), with
  Corollary~\ref{cor:caught}(i) for the clock or Corollary~\ref{cor:caught-threshold}(i) for the
  threshold, give $R=\mathcal C_{P,r_+}$ and the displayed nesting at the chosen margin, and
  $0<\Prb(B)<\Prb(A)$.
\item \emph{Pinned differentiable terms.} Condition (S8) pins $h^-$ and $h^+$. The maps
  $\kappa\mapsto\Eop[h^-\mid B]$ and
  $\kappa\mapsto\Eop[h^- -h^+\mid R]$ are differentiable. Conditions (S1) to (S3), (S7),
  (S10), and (S11) make the masses of $A$, $B$, and $R$ independent of $\kappa$.
\end{enumerate}
Define the caught-only leg, the re-pricing term oriented with the net cut, and the net cut leg by
\begin{align}
\label{eq:cut-split-definitions}
\widetilde s_B&=\partial_\kappa\Eop[h^-\mid B],
&\delta_{\mathrm{cut}}&=\partial_\kappa\Eop[h^- -h^+\mid R],
\nonumber\\
s_B&=\frac{1}{\Prb(B)}\partial_\kappa
\left(\Eop[h^-\ind_A]-\Eop[h^+\ind_R]\right).
\end{align}
Thus $\delta_{\mathrm{cut}}=-\delta$ for the clock notation above and
$\delta_{\mathrm{cut}}=-\delta_\tau$ for the threshold notation above. This orientation makes both
contributions enter the net cut with positive coefficients.
Then, at either margin,
\begin{equation}
\label{eq:cut-split-both}
s_B=\widetilde s_B+\frac{1-\varphi}{\varphi}\,\delta_{\mathrm{cut}}.
\end{equation}
The caught-only leg prices the newly flagged histories under the looser rule. The re-pricing term
prices the survivors under both rules.
\end{lemma}
\csname unlabelledfalse\endcsname

\begin{proof}
The sets $B$ and $R$ partition $A$, so
\[
\begin{aligned}
\Eop[h^-\ind_A]-\Eop[h^+\ind_R]
&=\Eop[h^-\ind_B]+\Eop[(h^- -h^+)\ind_R]\\
&=\Prb(B)\Eop[h^-\mid B]
  +\Prb(R)\Eop[h^- -h^+\mid R].
\end{aligned}
\]
The three masses are $\kappa$-free under (K-2). Differentiate, divide by $\Prb(B)>0$, and use
$\Prb(R)/\Prb(B)=(1-\varphi)/\varphi$. This gives
\eqref{eq:cut-split-both}.
\end{proof}

\csname unlabelledtrue\endcsname
\begin{lemma}[One crossing decides the composition band]
\label{lem:one-crossing}
Take a non-null nested clock or threshold cut at fixed policies and an interior liquidity
$\kappa\in(0,1)$ in the order-size-two model. Let $r_-$ be the looser rule and
$r_+$ the tighter rule, and write
\[
A=\mathcal C_{P,r_-},
\qquad
B=\mathcal C_{F,r_+}\setminus\mathcal C_{F,r_-},
\qquad
R=A\setminus B=\mathcal C_{P,r_+},
\qquad
\varphi=\frac{\Prb(B)}{\Prb(A)}\in(0,1).
\]
Assume standing conditions (S1) to (S13), so
Lemma~\ref{lem:g3} represents each pooled premium as a polynomial in $\kappa$ with coefficients that
do not depend on $\kappa$. In kernel units write
\begin{equation}
\label{eq:one-crossing-slopes}
s^q(\kappa)=-\frac12\sum_{k=0}^{H}(1-\varepsilon)^k\varepsilon^{H-k}c_k^q,
\qquad \varepsilon=\kappa/2,
\qquad q\in\{-,+\},
\end{equation}
where
$s^-=\partial_\kappa\Eop[h^-\mid A]$ and
$s^+=\partial_\kappa\Eop[h^+\mid R]$ are the signed pooled slopes under $r_-$ and $r_+$. Set
$d_k=c_k^+-c_k^-$ and, for $v=(v_0,\ldots,v_H)$, define the reversed polynomial
\begin{equation}
\label{eq:one-crossing-polynomial}
P_v(x)=\sum_{k=0}^{H}v_kx^{H-k},
\qquad x=x(\kappa)=\frac{\kappa}{2-\kappa}.
\end{equation}
Assume:
\begin{enumerate}[label=(R\arabic*),leftmargin=3.0em,itemsep=3pt]
\item \emph{One crossing for each pooled cell.} After zero coefficients are deleted, each of
  $(c_0^-,\ldots,c_H^-)$ and $(c_0^+,\ldots,c_H^+)$ changes sign exactly once, from negative to
  positive. The actual endpoint coefficients satisfy $c_0^-<0<c_H^-$ and $c_0^+<0<c_H^+$.
\item \emph{One crossing for the survivor difference.} After zero coefficients are deleted,
  $(d_0,\ldots,d_H)$ changes sign exactly once, from positive to negative, and
  $d_0>0>d_H$.
\item \emph{Liquidity is above the three revelation roots.} Let $r_A$, $r_R$, and $r_D$ be the
  positive roots of $P_{c^-}$, $P_{c^+}$, and $P_d$, which (R1), (R2), and Descartes' rule make
  unique. At the liquidity compared,
  \[
  x(\kappa)>r_*:=\max\{r_A,r_R,r_D\}.
  \]
\end{enumerate}
Then
\begin{equation}
\label{eq:one-crossing-order}
0<s^+(\kappa)<s^-(\kappa),
\end{equation}
and the net cut leg defined by
$s^-=(1-\varphi)s^++\varphi s_B$ satisfies
\begin{equation}
\label{eq:one-crossing-band}
s^-<s_B<\frac{s^-}{\varphi}
<\frac{2-\varphi}{\varphi}s^-.
\end{equation}
Thus the net cut leg lies in the composition band and the tighter pool has strictly lower liquidity
sensitivity. The composition ratio at the chosen margin, as defined in
Corollary~\ref{cor:caught} or Corollary~\ref{cor:caught-threshold}, is
$C_{\mathrm{cut}}=|s^+|/|s^-|<1$.
\end{lemma}
\csname unlabelledfalse\endcsname

\begin{proof}
The coefficient list of each $P_{c^q}$ has exactly one sign change. Descartes' rule of signs gives
exactly one positive root, counted with multiplicity. The strict endpoint coefficients give
$P_{c^q}(0)=c_H^q>0$ and a negative leading coefficient, so
$P_{c^q}(x)<0$ above its positive root. The same argument applied to $P_d$ gives one positive root,
with $P_d(0)=d_H<0$ and a positive leading coefficient, so $P_d(x)>0$ above that root.

\begin{equation}
\label{eq:one-crossing-factor}
\sum_{k=0}^{H}(1-\varepsilon)^k\varepsilon^{H-k}v_k
=(1-\varepsilon)^H P_v(x(\kappa)).
\end{equation}
The factor on the right is positive. Since $P_d=P_{c^+}-P_{c^-}$ by linearity of $v\mapsto P_v$,
condition (R3) therefore gives
\[
\Wrev^-<\Wrev^+<0,
\qquad
\Wrev^q:=\sum_{k=0}^{H}(1-\varepsilon)^k\varepsilon^{H-k}c_k^q.
\]
Equation~\eqref{eq:one-crossing-slopes} proves
\eqref{eq:one-crossing-order}. Solving the cut identity for the net cut leg gives
\[
s_B=\frac{s^--(1-\varphi)s^+}{\varphi}.
\]
Hence
\[
s_B-s^-=\frac{1-\varphi}{\varphi}(s^--s^+)>0,
\qquad
\frac{s^-}{\varphi}-s_B=\frac{1-\varphi}{\varphi}s^+>0.
\]
Since $0<\varphi<1$ and $s^->0$,
$s^-/\varphi<((2-\varphi)/\varphi)s^-$. This proves the band and the composition conclusion.
\end{proof}

\csname unlabelledtrue\endcsname
\begin{lemma}[Distinct crossings reverse total sensitivity]
\label{lem:reversal}
Compare a non-null nested pair at fixed policies, with $r_-$ the looser rule and $r_+$ the tighter
rule. Let $s^-$ be the signed pooled slope under $r_-$ and let $s^+$ be the slope under $r_+$,
using its own kernel on $\mathcal C_{P,r_+}$. Write
$\Omega_q=\Prb(\mathcal C_{F,r_q})$ for $q\in\{-,+\}$. Assume:
\begin{enumerate}[label=(V-\arabic*),leftmargin=3.2em,itemsep=3pt]
\item \emph{Standing conditions and factorisation.} Conditions (S1) to (S13) hold,
  $0<\Omega_-<1$ and $0<\Omega_+<1$, and the hypotheses of
  Proposition~\ref{prop:factorisation} hold under both rules on an open set in $(0,1)$ containing
  every zero of the two signed pooled slopes.
\item \emph{Distinct one-crossing roots.} The functions $s^-$ and $s^+$ are continuous on $(0,1)$.
  Each has exactly one zero there and changes sign at it. Write the roots as $\kappa_-$ and
  $\kappa_+$ respectively, and assume $\kappa_-\ne\kappa_+$.
\end{enumerate}
Then there is an open interval $I$ containing the looser pool's root $\kappa_-$ on which the tighter
rule is strictly more sensitive in total:
\begin{equation}
\label{eq:reversal}
\mathcal S(\kappa,r_+)>\mathcal S(\kappa,r_-)
\qquad\text{for every }\kappa\in I.
\end{equation}
\end{lemma}
\csname unlabelledfalse\endcsname

\begin{proof}
Factorisation, (S9), and Step~3 of Proposition~\ref{prop:factorisation}, which makes $\Omega_q$
constant in $\kappa$, give, for $q\in\{-,+\}$,
\[
\mathcal S(\kappa,r_q)=(1-\Omega_q)\Delta_m|s^q(\kappa)|.
\]
At the looser pool's root, $s^-(\kappa_-)=0$. The roots are distinct and $s^+$ has no other zero,
so $s^+(\kappa_-)\ne0$. Since $1-\Omega_+>0$ and $\Delta_m>0$,
\[
\mathcal S(\kappa_-,r_+)>0=\mathcal S(\kappa_-,r_-).
\]
Both sides are continuous in $\kappa$. Their strict difference at $\kappa_-$ therefore remains
positive on some open interval containing $\kappa_-$ and contained in the neighbourhood from
(V-1). This proves \eqref{eq:reversal}.
\end{proof}
